\documentclass[tikz,border=10pt]{standalone}
\usepackage{fontspec}
\usepackage[mathrm=sym]{unicode-math}
\setmainfont{texgyrepagella-regular.otf}[
    BoldFont = texgyrepagella-bold.otf,
    ItalicFont = texgyrepagella-italic.otf,
    BoldItalicFont = texgyrepagella-bolditalic.otf
]
\setmathfont{texgyrepagella-math.otf}[AutoFakeBold = 2.5]
\usepackage{tikz}
\usepackage{pgfplots}
\pgfplotsset{compat=1.18}
\usetikzlibrary{arrows.meta, calc}

\begin{document}
\begin{tikzpicture}[>=Stealth, line cap=round, line join=round]

  % Scale factors: q in fm (0 to 8.0), Delta in MeV (0 to 1.55)
  \def\xscale{1.10}
  \def\yscale{3.6}

  % Outer bounding frame
  \draw[line width=0.95pt] (0, 0) rectangle (8.0*\xscale, 1.55*\yscale);

  % Y-axis ticks and labels (left)
  \foreach \y/\lbl in {0.0/0.0, 0.5/0.5, 1.0/1.0} {
    \draw[line width=0.65pt] (0, \y*\yscale) -- (0.16, \y*\yscale);
    \node[anchor=east, font=\small] at (-0.08, \y*\yscale) {$\lbl$};
  }
  \node[anchor=south east, font=\normalsize, align=center] at (-0.15, 1.30*\yscale) {$\Delta(q, p_F(q))$\\[2pt]\small$(\mathrm{MeV})$};

  % X-axis ticks and labels (bottom)
  \foreach \q in {0, 1, 2, 3, 4, 5, 6, 7} {
    \draw[line width=0.65pt] (\q*\xscale, 0) -- (\q*\xscale, 0.16);
    \node[anchor=north, font=\small] at (\q*\xscale, -0.08) {$\q$};
  }
  \node[anchor=north west, font=\normalsize] at (7.15*\xscale, -0.08) {$q\,(\mathrm{fm})$};

  % ========================================================
  % UPPER CURVE: V0 = -28 MeV, r0 = 1.6 fm
  % ========================================================
  \draw[line width=1.15pt]
    plot[smooth] coordinates {
      (0*\xscale, 0.77*\yscale)
      (2.0*\xscale, 0.77*\yscale)
      (3.5*\xscale, 0.77*\yscale)
      (4.5*\xscale, 0.78*\yscale)
      (5.2*\xscale, 0.84*\yscale)
      (5.8*\xscale, 1.02*\yscale)
      (6.4*\xscale, 1.32*\yscale)
      (6.7*\xscale, 1.35*\yscale)
      (7.0*\xscale, 1.25*\yscale)
      (7.3*\xscale, 0.80*\yscale)
      (7.55*\xscale, 0.20*\yscale)
      (7.62*\xscale, 0.00*\yscale)
    };

  % Annotation for upper curve
  \node[anchor=south west, font=\small, align=left] (lblUp) at (3.5*\xscale, 1.15*\yscale)
    {$V_0 = -28\,\mathrm{MeV}$\\[1pt]$r_0 = 1.6\,\mathrm{fm}$};
  \draw[->, line width=0.75pt] (lblUp.south west) to[out=-135, in=65] (2.7*\xscale, 0.78*\yscale);

  % ========================================================
  % LOWER CURVE: V0 = 14 MeV, r0 = 2.7 fm
  % ========================================================
  \draw[line width=1.15pt]
    plot[smooth] coordinates {
      (0*\xscale, 0.60*\yscale)
      (3.0*\xscale, 0.60*\yscale)
      (5.0*\xscale, 0.61*\yscale)
      (5.8*\xscale, 0.63*\yscale)
      (6.4*\xscale, 0.72*\yscale)
      (6.8*\xscale, 0.95*\yscale)
      (7.15*\xscale, 1.31*\yscale)
      (7.30*\xscale, 1.32*\yscale)
      (7.45*\xscale, 1.10*\yscale)
      (7.60*\xscale, 0.50*\yscale)
      (7.68*\xscale, 0.00*\yscale)
    };

  % Annotation for lower curve
  \node[anchor=north west, font=\small, align=left] (lblDown) at (3.5*\xscale, 0.45*\yscale)
    {$V_0 = 14\,\mathrm{MeV}$\\[1pt]$r_0 = 2.7\,\mathrm{fm}$};
  \draw[->, line width=0.75pt] (lblDown.north west) to[out=135, in=-65] (2.7*\xscale, 0.59*\yscale);

\end{tikzpicture}
\end{document}
